% ++++++++++++++++++++++++++++++++++++++++++++++++++++++++++++++++++++++++
% Autor              : Ricardo Pereira Dias <contato@ricardopdias.com.br>
% Data de criação    : 29 de Janeiro de 2019
% Referência         : https://github.com/abntex/abntex2/blob/master/tex/latex/abntex2/abntex2.cls
% ++++++++++++++++++++++++++++++++++++++++++++++++++++++++++++++++++++++++

% Opções Memoir
% 12pt       tamanho da fonte
% openright  capítulos começam em pág ímpar (insere página vazia caso preciso)
% twoside    para impressão em recto e verso. Oposto a oneside
% a4paper    % tamanho do papel.

% Opções Abntex2
% Para converter as seções pra etra maiúscula
% chapter=TITLE, section=TITLE, subsection=TITLE, subsubsection=TITLE
% Para converter as seções pra letra minúscula
% chapter=title, section=title, subsection=title, subsubsection=title
% sumario=tradicional|abnt-6027-2012
% Opções do pacote Babel
% english   idioma adicional para hifenização
% french    idioma adicional para hifenização
% spanish   idioma adicional para hifenização
% brazil    % o último idioma é o principal do documento

\documentclass[
    12pt,openright,twoside,a4paper, % Memoir
    sumario=tradicional
    brazil, % Babel
]{abntex2}

% Pacotes configurados para utilização
% ++++++++++++++++++++++++++++++++++++++++++++++++++++++++++++++++++++++++
% Autor              : Ricardo Pereira Dias <contato@ricardopdias.com.br>
% Data de criação    : 29 de Janeiro de 2019
% ++++++++++++++++++++++++++++++++++++++++++++++++++++++++++++++++++++++++

% =========================================================================
% PACOTES
% =========================================================================


\usepackage[utf8]{inputenc}  % Codificacao do documento (conversão automática dos acentos)
\usepackage{indentfirst}     % Indenta o primeiro parágrafo de cada seção.
\usepackage{color}           % Controle das cores
\usepackage{graphicx}        % Inclusão de gráficos
\usepackage{microtype}       % para melhorias de justificação
\usepackage{xcolor}
\usepackage{fancyhdr}
\usepackage{longtable}
\usepackage{colortbl}
\usepackage{lipsum}       % para geração de dummy text
\usepackage[subentrycounter,seeautonumberlist,nonumberlist=true]{glossaries}
\usepackage[brazilian,hyperpageref]{backref}     % Paginas com as citações na bibl
\usepackage[alf]{abntex2cite}    % Citações padrão ABNT


% Informações pré-definidas
%=========================================================================
% Autor              : Ricardo Pereira Dias <rpdesignerfly@gmail.com>
%-------------------------------------------------------------------------
% Data de criação    : 29 de Janeiro de 2019
% Cores: https://www.overleaf.com/learn/latex/Using_colours_in_LaTeX
%=========================================================================

%=========================================================================
% CAPA e FOLHA DE ROSTO
%=========================================================================

% Projeto de Pesquisa
% TCC
% Artigo
\newcommand{\schoolInstitution}{%
    Universidade de Raccoon City - URC
    \par
    Faculdade de Arquitetura da Informação
    \par
    Programa de Pós-Graduação
}
\newcommand{\projectType}{Trabalho Acadêmico (TCC)}
\newcommand{\projectName}{Modelo Canônico de \projectType{} com Speed Latex}
\newcommand{\projectCity}{São Carlos - SP}
\newcommand{\projectDate}{2019}
\newcommand{\projectVersion}{, v-1.0.0}
% O preambulo deve conter o tipo do trabalho, o objetivo,
% o nome da instituição e a área de concentração
\newcommand{\projectPreamble}{Modelo canônico de  \projectType{} baseado no Abntex2 em conformidade com as normas ABNT, apresentado à comunidade de usuários \LaTeX.}
\newcommand{\projectAuthor}{Ricardo Pereira Dias}
\newcommand{\projectContact}{contato@ricardopdias.com.br}

% TCC
\newcommand{\projectOrientingLabel}{Prof.}%rotulo
\newcommand{\projectOrienting}{Albert Wesker}
\newcommand{\projectCoorientingLabel}{Prof.}%rotulo
\newcommand{\projectCoorienting}{Barry Burton}

% Artigo
\newcommand{\projectForeignName}{Canonical Academic Work template in SpeedLatex} % Nome extrangeiro

%=========================================================================
% CABEÇALHO
%=========================================================================
% Aparência
\definecolor{header_rule}{RGB}{ 0, 0, 0 }
\definecolor{header_text}{RGB}{ 0, 0, 0 }
\newcommand{\headerRuleHeight}{0pt}
% Páginas ímpares
\newcommand{\headerOddLeft}{}
\newcommand{\headerOddCenter}{}
\newcommand{\headerOddRight}{\thepage}
% Páginas Pares
\newcommand{\headerEvenLeft}{\thepage}
\newcommand{\headerEvenCenter}{}
\newcommand{\headerEvenRight}{}


%=========================================================================
% RODAPÉ
%=========================================================================
% Aparência
\definecolor{footer_rule}{RGB}{ 0, 0, 0 }
\definecolor{footer_text}{RGB}{ 0, 0, 0 }
\newcommand{\footerRuleHeight}{0pt}
% Páginas ímpares
\newcommand{\footerOddLeft}{}
\newcommand{\footerOddCenter}{}
\newcommand{\footerOddRight}{}
% Páginas Pares
\newcommand{\footerEvenLeft}{}
\newcommand{\footerEvenCenter}{}
\newcommand{\footerEvenRight}{}

%=========================================================================
% CORES GERAIS
%=========================================================================
% http://erikasarti.com/html/tabela-cores/
\definecolor{black}{RGB}{0,0,0}
% \definecolor{blue}{RGB}{41,5,195}
\definecolor{blue}{RGB}{0,0,0}
\definecolor{dim_gray}{RGB}{105,105,105}
\definecolor{gray}{RGB}{128,128,128}
\definecolor{dark_gray}{RGB}{169,169,169}
\definecolor{silver}{RGB}{192,192,192}
\definecolor{light_grey}{RGB}{211,211,211}
\definecolor{gainsboro}{RGB}{220,220,220}
\definecolor{white}{RGB}{255,255,255}

\definecolor{success_text}{RGB}{ 43, 180, 54 }
\definecolor{info_text}{RGB}{ 2, 69, 156 }
\definecolor{warning_text}{RGB}{ 255, 158, 0}
\definecolor{danger_text}{RGB}{ 255, 60, 0 }

%=========================================================================
% CAPÍTULOS E SEÇÕES
%=========================================================================

% Vários exemplos de capitulos
% http://tug.ctan.org/info/MemoirChapStyles/
% https://texblog.org/2012/07/03/fancy-latex-chapter-styles/

% Para formatar todos os niveis de seção
\newcommand{\sectioningTopStyle}{\bfseries}
\newcommand{\sectioningSecStyle}{\bfseries}
\newcommand{\sectioningFinalDot}{} % Para poder adicionar um character no final da numeração


% Partes
% Declaração: Parte 1
\newcommand{\partAssertionFont}{\sffamily}
\newcommand{\partAssertionStyle}{\sectioningTopStyle}
\newcommand{\partAssertionSize}{\large}
% Título da parte
\newcommand{\partFont}{\sffamily}
\newcommand{\partStyle}{\sectioningTopStyle}
\newcommand{\partSize}{\Huge}
\definecolor{part_count}{RGB}{ 0, 0, 0 }
\definecolor{part_text}{RGB}{ 0, 0, 0 }

% Capítulos
\newcommand{\chapterFont}{\sffamily}
\newcommand{\chapterStyle}{\sectioningTopStyle}
\newcommand{\chapterSize}{\Huge}
\definecolor{chapter_count}{RGB}{ 0, 0, 0 }
\definecolor{chapter_text}{RGB}{ 0, 0, 0 }

% Seções
\newcommand{\sectionFont}{\sffamily}
\newcommand{\sectionStyle}{\sectioningSecStyle}
\newcommand{\sectionSize}{\Large}
\definecolor{section_count}{RGB}{ 0, 0, 0 }
\definecolor{section_text}{RGB}{ 0, 0, 0 }

% Sub Seções
\newcommand{\subsectionFont}{\sffamily}
\newcommand{\subsectionStyle}{\sectioningSecStyle}
\newcommand{\subsectionSize}{\large}
\definecolor{subsection_count}{RGB}{ 0, 0, 0 }
\definecolor{subsection_text}{RGB}{ 0, 0, 0 }

% Sub Sub Seções
\newcommand{\subsubsectionFont}{\sffamily}
\newcommand{\subsubsectionStyle}{\sectioningSecStyle}
\newcommand{\subsubsectionSize}{\normalsize}
\definecolor{subsubsection_count}{RGB}{ 0, 0, 0 }
\definecolor{subsubsection_text}{RGB}{ 0, 0, 0 }

% Sub Sub Sub Seções
\newcommand{\subsubsubsectionFont}{\sffamily}
\newcommand{\subsubsubsectionStyle}{\sectioningSecStyle}
\newcommand{\subsubsubsectionSize}{\normalsize}
\definecolor{subsubsubsection_count}{RGB}{ 0, 0, 0 }
\definecolor{subsubsubsection_text}{RGB}{ 0, 0, 0 }

\definecolor{paragraph_text}{RGB}{ 0, 0, 0 } % parágrafo de leis ou documentos
\definecolor{subparagraph_text}{RGB}{ 0, 0, 0 } % subparágrafo de leis ou documentos

% Corpo de texto
\definecolor{body_text}{RGB}{ 0, 0, 0 }
\definecolor{link_text}{RGB}{ 0, 0, 0 }


% Configurações do documento
\include{libraries/document-config-common}

% Estilos de capítulos
% =========================================================================
% Autor              : Ricardo Pereira Dias <rpdesignerfly@gmail.com>
% Data de criação    : 30 de Março de 2019
% -------------------------------------------------------------------------

%% =========================================================================
% Autor              : Ricardo Pereira Dias <rpdesignerfly@gmail.com>
% Data de criação    : 30 de Março de 2019
% -------------------------------------------------------------------------

% Formatação das seções
\usepackage{titlesec}

% Declaração dos capítulos
\titleformat{\chapter}[display]
    {\normalfont\huge\bfseries}{\color{chapter_label}\chaptertitlename\  \color{chapter_count}\thechapter}{20pt}{\Huge}


% Fontes que acompanham o texlive
% % =========================================================================
% Autor              : Ricardo Pereira Dias <rpdesignerfly@gmail.com>
% Data de criação    : 30 de Março de 2019
% =========================================================================
% Para listar, no terminal: "fc-list : family"
% -------------------------------------------------------------------------
% Links interssantes:
% http://www.tug.dk/FontCatalogue/
% http://www.tug.dk/FontCatalogue/allfonts.html
% http://www.tug.dk/FontCatalogue/cabin/
% http://mirrors.ibiblio.org/CTAN/fonts/
% -------------------------------------------------------------------------

% code: pag

\usepackage{avant}
\usepackage[T1]{fontenc}

% % =========================================================================
% Autor              : Ricardo Pereira Dias <rpdesignerfly@gmail.com>
% Data de criação    : 30 de Março de 2019
% =========================================================================
% Para listar todas as fontes disponiveis no sistema,
% digite no terminal:
% fc-list : family
% -------------------------------------------------------------------------
% Links interssantes:
% http://www.tug.dk/FontCatalogue/
% https://en.wikibooks.org/wiki/LaTeX/Fonts
% https://pt.overleaf.com/learn/latex/Font_sizes,_families,_and_styles#Font_families
% -------------------------------------------------------------------------

% https://ctan.org/tex-archive/fonts/urw/arial

\usepackage{uarial}
\usepackage[T1]{fontenc}

% % =========================================================================
% Autor              : Ricardo Pereira Dias <rpdesignerfly@gmail.com>
% Data de criação    : 30 de Março de 2019
% =========================================================================
% OpenSans
% -------------------------------------------------------------------------
% Para listar todas as fontes disponiveis no sistema,
% fc-list : family
% -------------------------------------------------------------------------
% Links interssantes:
% http://www.tug.dk/FontCatalogue/
% http://www.tug.dk/FontCatalogue/allfonts.html
% http://www.tug.dk/FontCatalogue/cabin/
% -------------------------------------------------------------------------

% https://ctan.org/tex-archive/fonts/bera

% Bera Serif: fve
% Bera Sans: fvs
% Bera Sans Mono: fvm

\usepackage{bera}
\usepackage[T1]{fontenc}

% \include{libraries/fonts/cabin}
% % =========================================================================
% Autor              : Ricardo Pereira Dias <rpdesignerfly@gmail.com>
% Data de criação    : 30 de Março de 2019
% =========================================================================
% Para listar todas as fontes disponiveis no sistema,
% digite no terminal:
% fc-list : family
% -------------------------------------------------------------------------
% Links interssantes:
% http://www.tug.dk/FontCatalogue/
% http://www.tug.dk/FontCatalogue/firasansbook/
% -------------------------------------------------------------------------

% https://ctan.org/tex-archive/fonts/fira

\usepackage[book]{FiraSans}
\usepackage[T1]{fontenc}
\renewcommand*\oldstylenums[1]{{\firaoldstyle #1}}

% % =========================================================================
% Autor              : Ricardo Pereira Dias <rpdesignerfly@gmail.com>
% Data de criação    : 29 de Janeiro de 2019
% =========================================================================
% COMPUTER MODERN
% -------------------------------------------------------------------------
% Para listar todas as fontes disponiveis no sistema,
% digite no terminal:
% fc-list : family
% -------------------------------------------------------------------------
% Links interssantes:
% http://www.tug.dk/FontCatalogue/
% https://en.wikibooks.org/wiki/LaTeX/Fonts
% https://pt.overleaf.com/learn/latex/Font_sizes,_families,_and_styles#Font_families
% -------------------------------------------------------------------------

\usepackage{lmodern}
\usepackage[T1]{fontenc}

\include{libraries/fonts/roboto}

% Inclui as entradas do glossario
\include{templates/common-book/glossary}

\begin{document}

    % Define o idioma do documento (conforme pacotes do babel)
    \selectlanguage{brazil}

    % Retira espaço extra obsoleto entre as frases.
    \frenchspacing

    %---
    %
    % Se desejar escrever o artigo em duas colunas, descomente a linha abaixo
    % e a linha com o texto ``FIM DE ARTIGO EM DUAS COLUNAS''.
    % \twocolumn[    		% INICIO DE ARTIGO EM DUAS COLUNAS
    %
    %---

    %===============================================
    % ELEMENTOS PRÉ-TEXTUAIS
    %===============================================
    %% ----------------------------------------------
% Capa
% ----------------------------------------------

\begin{capa}%

    \center
    
    \includegraphics[scale=.5]{assets/logo.eps}

    \vfill

    \chapterFont\bfseries\LARGE\projectName

    \vfill

    \chapterFont\large\projectAuthor

    \vfill

    \large\projectCity
    \large\projectDate

    \vspace*{1cm}
\end{capa}
%
    % %% ----------------------------------------------
% Changelog
% ----------------------------------------------

    \def\arraystretch{1.5}%  1 is the default, change whatever you need

    \footnotesize
    \begin{longtable}{%
        |m{2cm}|m{5.5cm}|m{3cm}|m{3cm}|
    }%

        \hline

        \rowcolor{gray}\textcolor{white}{Data}
        & \textcolor{white}{Descrição}
        & \textcolor{white}{Autor}
        & \textcolor{white}{Revisor} \\

        \hline

        \small{10/01/1980}
        & \small{\lipsum[2]}
        & \small{Ricardo Pereira}
        & \small{Tiryon Lenister} \\

        \hline

        \small{10/01/1980}
        & \small{\lipsum[2]}
        & \small{Ricardo Pereira}
        & \small{Tiryon Lenister} \\

        \hline

        \small{10/01/1980}
        & \small{\lipsum[2]}
        & \small{Ricardo Pereira}
        & \small{Tiryon Lenister} \\

        \hline

    \end{longtable}

\lipsum[2]

\vfill

\pagebreak
%
    \input{templates/common-book/pre-textuais/03-sumario}%


    %===============================================
    % ELEMENTOS TEXTUAIS
    %===============================================
    \textual

    \part{Sobre o Projeto}

    
% ----------------------------------------------------------
% Introdução (exemplo de capítulo sem numeração, mas presente no Sumário)
% ----------------------------------------------------------
\chapter{Introdução}
% ----------------------------------------------------------

Este documento e seu código-fonte são exemplos de referência de uso do Speed Latex,
que faz uso da excelente classe \textsf{memoir} e da customização oferecida pelo projeto \textsf{\abnTeX}. Optou-se por utilizar a classe \textsf{\abnTeX} sem alterá-la, seguindo a ``recomendação dos autores do projeto''\footnote{\url{https://github.com/abntex/abntex2/wiki/ComoCustomizar}}.

\textbf{Atenção:} Embora faça-se uso do \textsf{\abnTeX}, o modelo "Livro Simples"\ não segue as regras da ABNT.

\section{Modelos disponíveis}

O Speed Latex disponibiliza vários modelos prontos para criação de documentos baseados em Latex.
Confira abaixo a lista de modelos disponíveis:

\begin{itemize}

    \item \textbf{abnt-academic-work} - Para Trabalhos Acadêmicos/TCC (ABNT);
    \item \textbf{abnt-article} - Para Artigos Científicos (ABNT);
    \item \textbf{abnt-research-project} - Para Projetos de Pesquisa (ABNT);
    \item \textbf{common-letter} - Para criação de Cartas;
    \item \textbf{common-document} - Para criação de Documentos;
    \item \textbf{common-book} - Para criação de Livros

\end{itemize}

\section{Organização}

Além de possibilitar a criação de projetos pré-organizados baseados em Latex, o Speed Latex
permite uma maior facilidade na personalização do documento. Isso é feito através de um arquivo de configuração existente para cada projeto, onde possível personalizar várias informações, como:

\begin{itemize}

    \item O conteúdo dos cabeçalhos e rodapés;
    \item As cores das seções e dos textos em geral;
    \item Definir informações padrões para uso no documento

\end{itemize}

\section{Atualizações}

O Speed Latex \footnote{\url{https://github.com/ricardopedias/speed-latex}} está em constante atualização. Sinta-se à vontade para conferir os releases \footnote{\url{https://github.com/ricardopedias/speed-latex/releases}} e acompanhar as atualizações.

Sinceramente espera-se que cada atualização do Speed Latex aprimore a qualidade do trabalho que você produzirá, de modo que o principal esforço seja concentrado no conteúdo produzido e não na formatação de documentos.
%

    \part{Personalizando}

    
% ----------------------------------------------------------
% Introdução (exemplo de capítulo sem numeração, mas presente no Sumário)
% ----------------------------------------------------------
\chapter{Cabeçalhos e Rodapés}

Para personalizar os cabeçalhos, basta definir as informações correspondentes no arquivo de configuração, de forma que se adequem à aparência desejada:

\begin{itemize}

    \item abnt-academic-work - Para Trabalhos Acadêmicos/TCC (ABNT);

\end{itemize}

\chapter{Naitis e Rodapés}

Para personalizar os cabeçalhos, basta definir as informações correspondentes no arquivo de configuração, de forma que se adequem à aparência desejada:

\begin{itemize}

    \item abnt-academic-work - Para Trabalhos Acadêmicos/TCC (ABNT);

\end{itemize}
%
    

\chapter{Tabelas}

\lipsum[2]

\def\arraystretch{1.3}%  1 is the default, change whatever you need

% p{width}	Top-aligned column
% m{width}	Middle-aligned column
% b{width}	Bottom-aligned column
% >{\raggedright}	Left-aligned column
% >{\centering}	Center-aligned column
% >{\raggedleft}	Right-aligned column

% \usepackage{array}
%
% % Any column
% >{\centering}p{0.1\linewidth}
%
% % Last column
% >{\centering\arraybackslash}p{0.1\linewidth}

\footnotesize
\begin{longtable}{%
    |
    >{\centering}m{2cm}|
    >{\raggedright}m{5.5cm}|
    >{\centering}m{3cm}|
    >{\centering\arraybackslash}m{3cm}
    |
}%

    \hline

    \rowcolor{gray}\textcolor{white}{Data}
    & \textcolor{white}{Descrição}
    & \textcolor{white}{Autor}
    & \textcolor{white}{Revisor} \\

    \hline

    \small{10/01/1980}
    & \small{\lipsum[2]}
    & \small{Ricardo Pereira}
    & \small{Tiryon Lenister} \\

    \hline

    \small{10/01/1980}
    & \small{\lipsum[2]}
    & \small{Ricardo Pereira}
    & \small{Tiryon Lenister} \\

    \hline

    \small{10/01/1980}
    & \small{\lipsum[2]}
    & \small{Ricardo Pereira}
    & \small{Tiryon Lenister} \\

    \hline

\end{longtable}

\lipsum[2-8]
%

    \part{Usando Latex}

    
% ---------------------------
\chapter{Lorem ipsum dolor sit amet}

\section{Aliquam vestibulum}
\lipsum[2]

\section{Fringilla lorem}
\lipsum[3-5]

\section{Pellentesque habitant}
\lipsum[6-7]

% ---------------------------
\chapter{Donec varius orci}

\section{Cras viverra metus}
\lipsum[8-9]

\section{Nulla malesuada porttitor}
\lipsum[10-11]
%

    % ---
    % Finaliza a parte no bookmark do PDF, para que se inicie o bookmark na raiz
    % ---
    \bookmarksetup{startatroot}%

    % 
    %===================================================================================================
    % ELEMENTOS TEXTUAIS
    %===================================================================================================

    % \pagestyle{meuestilo}
    % \textual

    % ------------------------------------------------------------------------
    % Conclusão
    % ------------------------------------------------------------------------
    \chapter*[Considerações finais]{Considerações finais}
    \addcontentsline{toc}{chapter}{Considerações finais}

    \lipsum[31-33]
%


    %===============================================
    % ELEMENTOS PÓS-TEXTUAIS
    %===============================================
    \postextual

    % ----------------------------------------------------------
    % Referências bibliográficas
    % ----------------------------------------------------------
    \bibliography{templates/common-book/pos-textuais/01-references}
    % \input{templates/common-book/pos-textuais/02-glossario}
    % \input{templates/common-book/pos-textuais/03-apendices}
    % \cftinserthook{toc}{AAA}
    % 
% ----------------------------------------------
% Anexos
% ----------------------------------------------
\begin{anexosenv}

    % ---
    \chapter{Cras non urna sed feugiat cum sociis natoque penatibus et magnis dis
    parturient montes nascetur ridiculus mus}
    % ---

    \lipsum[31]

    \vspace*{1.5cm}

\end{anexosenv}

    % \input{templates/common-book/pos-textuais/05-agradecimentos}

\end{document}
